\documentclass[10pt]{article}
\usepackage[a4paper,top=1.55cm,bottom=1.7cm,left=1.55cm,right=1.55cm,columnsep=0.55cm]{geometry}
\usepackage{fontspec}
\usepackage{amsmath,amssymb}
\usepackage{booktabs,array}
\usepackage{multicol}
\usepackage{algpseudocode}
\usepackage{microtype}
\usepackage[hidelinks]{hyperref}
\setmainfont{STSong}[BoldFont=Heiti SC]
\setsansfont{Arial}
\setmonofont{Menlo}
\XeTeXlinebreaklocale "zh"
\XeTeXlinebreakskip=0pt plus 1pt
\setlength{\parindent}{1.5em}
\setlength{\parskip}{0pt}
\setlength{\abovedisplayskip}{4pt}
\setlength{\belowdisplayskip}{4pt}
\setlength{\textfloatsep}{7pt}
\setlength{\floatsep}{6pt}
\pagestyle{plain}
\newcommand{\code}[1]{\texttt{#1}}
\newcounter{algo}
\begin{document}

\setcounter{section}{2}
\section{方法}

\begin{center}
\fbox{\begin{minipage}{0.94\textwidth}\centering
\textbf{可追溯风险抽象：}
L0 原始经历 $\mid$ L1 风险记忆卡 $\mid$ L2 风险规则\\[3pt]
\textbf{面向职责的证据分配：}
Task Advocate $\mid$ Risk Critic $\mid$ Safety Decision Agent\\[3pt]
\textbf{执行权限边界：}
结构化建议 $\mid$ 确定性约束优化器 $\mid$ Safety Guard
\end{minipage}}
\refstepcounter{figure}\label{fig:overview}
\vspace{2pt}
{\small 图~\thefigure：方法概览。框架构造可追溯的风险证据，按决策职责分配证据，并分离语义建议与动作执行。}
\end{center}

\begin{multicols}{2}

\subsection{系统概览}

本文提出一种风险经验驱动的多角色框架，其核心包含两个相互衔接的操作：\textbf{可追溯风险抽象}与\textbf{面向职责的证据分配}。给定交互历史 $\mathcal{H}_t$，前者逐步生成单案例记忆卡和跨案例风险规则，并保留它们与原始交互记录的关联。给定当前规划输入 $\mathcal{X}_t=(g_t,r_t,e_t,\mathcal{A}_t)$，后者为每个决策角色构造不同的证据包。二者结合，使系统能够使用紧凑的长期风险证据，同时保留其经验来源。

Task Advocate、Risk Critic 和 Safety Decision Agent 分别解释所分配的证据，并输出包含证据 ID 的结构化建议，不直接产生可执行控制。确定性约束优化器从允许动作集 $\mathcal{A}_t$ 中计算最终语义动作，独立 Safety Guard 在执行前对其进行验证。第 3.2 节说明可追溯风险证据的构造，第 3.3 节说明证据的角色条件化分配及其决策用途。

\subsection{可追溯的分层风险记忆}

本文通过逐层抽象构造风险证据。L0 经历集 $\mathcal{E}_t$ 保留完整交互，L1 卡片集 $\mathcal{M}_t$ 提取单个案例的风险结构，L2 规则集 $\mathcal{Q}_t$ 归纳由多个案例共同支持的风险模式。每层抽象均通过稳定 ID 关联其来源；这些关联用于校验决策引用、衡量规则的经验支持，并在新证据与规则冲突时定位需要重新评估的案例。因此，该层次结构在压缩决策上下文的同时，使抽象结果仍能由原始证据验证和更新。

\subsubsection{L0：具身风险经历表示}

L0 将一次交互表示为经历 $E_i$，记录任务与场景条件、动作 $A_i$、结果 $O_i$、失败原因、风险类型、严重度 $\sigma_i\in\{1,\ldots,5\}$ 和轨迹来源。不可变的 \code{experience\_id} 将其链接至传感日志、控制记录或人工标注。

为在统一的风险尺度下比较不同观测，本文将 $E_i$ 编码为风险状态 $z_i=[1-\bar c_i,1-\bar d_i,1-p_i,1-b_i,h_i]$。其中，$\bar c_i$ 和 $\bar d_i$ 是归一化通行余量与障碍物距离，$p_i$ 是观测置信度，$b_i$ 是电量比例，$h_i$ 包含机器人、障碍和任务等离散条件。取补值后，较大的前四项均表示更强的风险因素；完整轨迹保留在 L0 中供审计。

\subsubsection{L1：风险记忆卡}

L1 将每条经历压缩为风险记忆卡。归一化严重度 $\sigma_i/5$ 表示后果强度；结果权重 $\omega(O_i)$ 对 failure、aborted 和 success 分别取 $1$、$0.8$ 和 $0.05$；可靠性 $\gamma_i\in[0,1]$ 用于降低低质量记录的影响。三者乘积 $q_i=(\sigma_i/5)\omega(O_i)\gamma_i$ 是相对证据分数，而非碰撞概率。令 $F_i$ 表示失败原因，$B_i$ 表示缓解措施与停止条件，$P_i$ 表示由 memory ID、experience ID 和来源片段组成的来源记录。

将 $z_i$ 与 $q_i$ 逐步代入压缩映射，可得 L1 卡片：
\begin{equation}
\begin{aligned}
M_i&=C(E_i)=\operatorname{Pack}(z_i,A_i,O_i,F_i,B_i,q_i,P_i)\\
&=\operatorname{Pack}\!\left([1-\bar c_i,1-\bar d_i,1-p_i,1-b_i,h_i],\right.\\
&\hspace{3.6em}\left.A_i,O_i,F_i,B_i,
\frac{\sigma_i}{5}\omega(O_i)\gamma_i,P_i\right).
\end{aligned}
\label{eq:card-derivation-zh}
\end{equation}
式~\eqref{eq:card-derivation-zh} 表明，L0 经历的风险状态、结果证据和来源记录共同构成可追溯的 L1 卡片。

\subsubsection{L2：风险规则聚合与更新}

L2 将风险状态、风险类型和动作语义相近的卡片组成模式集合 $\mathcal{I}_j$。聚合映射 $\mathcal{A}$ 根据其中的触发条件、建议与禁止动作、缓解措施和来源 ID 生成规则 $Q_j$。令 $n_j^+$、$n_j^-$ 和 $n_j^0$ 分别表示缓解成功、风险失败和未决中止次数，并以 $\theta_j$ 表示新观测为规则 $Q_j$ 提供明确证据的概率。采用 $\operatorname{Beta}(\alpha,\beta)$ 先验，其后验均值 $c_j$ 为
\begin{equation}
\begin{aligned}
(Q_j,c_j)&=\mathcal{A}(\{M_i\}_{i\in\mathcal{I}_j}),\\
c_j&=\mathbb{E}[\theta_j\mid\mathcal{I}_j]
=\frac{n_j^++n_j^-+\alpha}
{n_j^++n_j^-+n_j^0+\alpha+\beta}.
\end{aligned}
\label{eq:rule-zh}
\end{equation}
这里 $c_j$ 衡量明确经验支持的程度，而非风险发生概率。$Q_j$ 的来源集合决定 $c_j$ 中的计数，并在规则转为 \textit{active} 前接受一致性检查。当新的 L0 结果支持或反驳该规则时，来源 ID 用于定位相关卡片，进而重新计算支持度、修订触发条件或将规则转为 \textit{retired}。因此，版本化的 $Q_j\rightarrow M_i\rightarrow E_i$ 链参与规则验证与维护，而非仅作为附加元数据保存。

运行时，规则闸门仅允许所有命中 \textit{active} 规则共同建议且均未禁止的动作；若交集为空，则停止执行并请求人工复核。

\subsection{角色条件化的风险证据检索}

同一风险历史服务于不同的决策职责，因此本文按角色分配证据，而非向所有角色提供同一检索结果。Advocate 接收与任务推进和成功缓解相关的证据，Critic 接收失败案例与停止条件，Decision Agent 接收用于权衡任务收益与风险的证据。给定 $\mathcal{X}_t$，每张卡片产生特征向量 $\phi_i=[\mathrm{sim},\mathrm{sev},\mathrm{rel},\mathrm{rec},\mathrm{fit}]\in[0,1]^5$，分别表示情境相似度、严重度、可靠性、时效性和角色适配度。角色权重 $w^{(r)}$ 将这些公共特征转化为面向职责的相关性。

对于角色 $r$，$\mathcal{B}$ 表示候选证据包，$K$ 限制卡片数量，$B_r$ 是 Token 预算。冗余度 $d(i,j)$ 用于减少风险类型、结果和缓解措施重复的案例，$\widehat{\tau}(M_i)$ 估计卡片 $M_i$ 的 Token 开销，$\lambda$ 控制冗余惩罚。检索据此选择
\begin{equation}
\begin{aligned}
\mathcal{B}_r^*=\arg\max_{\mathcal{B}\subseteq\mathcal{M}_t}\;&
\sum_{M_i\in\mathcal{B}}(w^{(r)})^\top\phi_i
-\lambda\!\sum_{\substack{M_i,M_j\in\mathcal{B}\\i<j}}d(i,j)\\
\text{s.t. }\;&|\mathcal{B}|\leq K,\quad
\sum_{M_i\in\mathcal{B}}\widehat{\tau}(M_i)\leq B_r .
\end{aligned}
\label{eq:retrieval-zh}
\end{equation}
对于 $(\mathrm{sim},\mathrm{sev},\mathrm{rel},\mathrm{rec},\mathrm{fit})$，本文取
\[
\begin{aligned}
w^{(\mathrm{adv})}&=(0.35,0.05,0.20,0.05,0.35),\\
w^{(\mathrm{crit})}&=(0.30,0.25,0.15,0.05,0.25),\\
w^{(\mathrm{dec})}&=(0.35,0.20,0.20,0.10,0.15).
\end{aligned}
\]
算法~\ref{alg:retrieval} 以贪心边际增益近似求解式~\eqref{eq:retrieval-zh}，实验部分评估分角色权重的作用。

每个证据包记录角色与检索版本、选中的 L1 卡片和命中的 L2 规则、得分、匹配原因、预算用量及来源 ID。命中规则约束可行动作，其来源集合则记录该约束由哪些案例支持。Agent 通过 \code{cited\_experience\_ids} 和 \code{cited\_rule\_ids} 声明所用证据，框架拒绝证据包外的 ID，并将合法引用与优化器及 Safety Guard 输出共同保存，从而把动作的执行或拒绝关联到相应规则及其来源案例。

令 $\mathcal{D}_t$ 表示三个角色的报告，$\mathcal{F}_t$ 表示满足硬约束和 L2 禁止规则的动作集合。对于动作 $a$，$U(a)$ 是任务效用，$R(a\mid\mathcal{B}_{\mathrm{dec}}^*)$ 是场景与记忆的联合风险，$Z(a\mid\mathcal{D}_t,\mathcal{B}_{\mathrm{dec}}^*)$ 是由观测置信度、记忆覆盖和角色分歧形成的不确定性。阈值 $\rho$ 和 $\eta$ 分别限制这两个相对分数。确定性优化器选择
\begin{equation}
\begin{aligned}
a_t^*&=\arg\max_{a\in\mathcal{F}_t}U(a),\\
R(a\mid\mathcal{B}_{\mathrm{dec}}^*)&\leq\rho,\\
Z(a\mid\mathcal{D}_t,\mathcal{B}_{\mathrm{dec}}^*)&\leq\eta.
\end{aligned}
\label{eq:decision-zh}
\end{equation}
规则冲突、输出无效、可行集为空或 Safety Guard 拒绝动作时，框架触发 \code{safe\_stop}、\code{observe\_again} 或人工复核。该执行边界允许 Agent 解释检索证据，同时由确定性模块保留动作权限。

\begin{center}
\fbox{\begin{minipage}{0.91\columnwidth}
\refstepcounter{algo}\label{alg:retrieval}
\textbf{算法~\thealgo：角色条件化风险检索}
\begingroup
\algrenewcommand\algorithmicrequire{\textbf{输入：}}
\algrenewcommand\algorithmicensure{\textbf{输出：}}
\algrenewcommand\algorithmicfor{\textbf{对}}
\algrenewcommand\algorithmicforall{\textbf{对所有}}
\algrenewcommand\algorithmicdo{\textbf{执行}}
\algrenewcommand\algorithmicwhile{\textbf{当}}
\algrenewcommand\algorithmicend{\textbf{结束}}
\algrenewcommand\algorithmicif{\textbf{若}}
\algrenewcommand\algorithmicthen{\textbf{则}}
\algrenewtext{EndFor}{\textbf{循环结束}}
\algrenewtext{EndWhile}{\textbf{循环结束}}
\algrenewtext{EndIf}{\textbf{条件结束}}
\begin{algorithmic}[1]
\footnotesize
\Require 状态 $\mathcal{X}_t$，角色集 $\mathcal{R}$，L1 卡片集 $\mathcal{M}$，active L2 规则集 $\mathcal{Q}$，预算 $\{B_r\}$，上限 $K$
\Ensure 分角色证据包 $\{\mathcal{B}_r\}_{r\in\mathcal{R}}$
\State 校验 $\mathcal{X}_t$、$\mathcal{R}$、$\mathcal{M}$ 与 $\mathcal{Q}$
\For{每个角色 $r\in\mathcal{R}$}
  \State 计算所有 $\phi_i$；初始化候选集 $\mathcal{C}_r$ 与 $\mathcal{B}_r\gets\varnothing$
  \While{$\mathcal{C}_r\neq\varnothing$ 且 $|\mathcal{B}_r|<K$}
    \State 按式~\eqref{eq:retrieval-zh} 的边际增益选取 $M_{i^*}$
    \State 从 $\mathcal{C}_r$ 删除 $M_{i^*}$
    \If{$\widehat{\tau}(\mathcal{B}_r)+\widehat{\tau}(M_{i^*})\leq B_r$}
      \State $\mathcal{B}_r\gets\mathcal{B}_r\cup\{M_{i^*}\}$
    \EndIf
  \EndWhile
  \State 向 $\mathcal{B}_r$ 附加命中的 active L2 规则、得分、原因和来源 ID
\EndFor
\State \Return $\{\mathcal{B}_r\}_{r\in\mathcal{R}}$
\end{algorithmic}
\endgroup
\end{minipage}}
\end{center}

\end{multicols}
\end{document}
